The Reock score provides a boundary-insensitive compactness measure that
complements Polsby-Popper in redistricting litigation.
Our principal findings are:

\begin{enumerate}
  \item Reock decreases by only $0.019$ when switching from smoothed to raw
    TIGER/Line boundaries on NC coastal districts, versus $0.048$ for PP---a
    2.5$\times$ advantage in boundary stability.
  \item The moving-knife algorithm achieves mean Reock $\geq 0.38^\dagger$ on
    NC 2020, a 23\% improvement over standard-bisect, validating MKA's
    Reock maximisation design objective.
  \item The empirical Reock--PP gap averages $+0.15^\dagger$ on NC 2020
    under ratio-optimal, confirming that Reock is systematically more lenient
    than PP as a compactness criterion.
\end{enumerate}

Practitioners facing coastal-boundary challenges should present Reock as the
primary area-based metric, supported by the centroid-plus-max-radius approximation
described in Section~3 and the boundary sensitivity analysis in Section~4.
The MBC visual overlay is effective for courtroom presentation.

The K.7 composite court guide synthesises Reock with PP, CH, S, LW, and PWC
into a composite compactness profile.
